%% ============================================================
%% Clustering-Guided Contrastive Learning for Clinically Aware
%% Chest X-ray Image--Report Retrieval
%% Author: Diep Duc Lai
%% Compile: pdflatex paper_en.tex (x2)
%% ============================================================
\documentclass[11pt, a4paper]{article}

% -------- Encoding --------
\usepackage[utf8]{inputenc}
\usepackage[T5]{fontenc}   % T5 supports Vietnamese diacritics
\usepackage{lmodern}

% -------- Maths --------
\usepackage{amsmath, amssymb, amsthm}
\usepackage{bm}

% -------- Layout --------
\usepackage[top=2.5cm, bottom=2.5cm, left=3.0cm, right=2.5cm]{geometry}
\usepackage{setspace}
\onehalfspacing
\usepackage{parskip}
\setlength{\parindent}{1.2em}

% -------- Tables & Figures --------
\usepackage{booktabs, multirow, tabularx, makecell}
\usepackage{graphicx}
\usepackage{caption, subcaption}
\usepackage{float}
\usepackage{array}
\newcolumntype{C}[1]{>{\centering\arraybackslash}m{#1}}
\newcolumntype{L}[1]{>{\raggedright\arraybackslash}m{#1}}
\newcolumntype{R}[1]{>{\raggedleft\arraybackslash}m{#1}}

% -------- Hyperref --------
\usepackage[colorlinks=true,
            linkcolor=blue!70!black,
            citecolor=green!50!black,
            urlcolor=blue!80!black]{hyperref}

% -------- References --------
\usepackage{natbib}
\bibliographystyle{abbrvnat}
\setcitestyle{authoryear,open={(},close={)}}

% -------- Misc --------
\usepackage{xcolor}
\usepackage{colortbl}
\usepackage{enumitem}
\usepackage{microtype}
\usepackage{fancyhdr}
\usepackage{titlesec}
\usepackage{lipsum}

\pagestyle{fancy}
\fancyhf{}
\fancyhead[L]{\small\textit{Clustering-Guided Contrastive Learning for Chest X-ray Image--Report Retrieval}}
\fancyhead[R]{\small\thepage}
\renewcommand{\headrulewidth}{0.4pt}
\fancypagestyle{plain}{\fancyhf{}\renewcommand{\headrulewidth}{0pt}}

% Section formatting
\titleformat{\section}{\large\bfseries}{\thesection.}{0.5em}{}
\titleformat{\subsection}{\normalsize\bfseries}{\thesubsection.}{0.5em}{}
\titleformat{\subsubsection}{\normalsize\itshape\bfseries}{\thesubsubsection.}{0.5em}{}

% -------- Custom commands --------
\newcommand{\calL}{\mathcal{L}}
\newcommand{\calQ}{\mathcal{Q}}
\newcommand{\vv}[1]{\mathbf{#1}}
\newcommand{\norm}[1]{\left\lVert#1\right\rVert}
\newcommand{\ie}{\textit{i.e.}}
\newcommand{\eg}{\textit{e.g.}}
\newcommand{\etal}{\textit{et al.}}

% =========================================================
\begin{document}
% =========================================================

%% ---------------------------------------------------------
%% TITLE PAGE
%% ---------------------------------------------------------
\begin{titlepage}
  \centering
  \vspace*{2cm}

  {\LARGE\bfseries
    Clustering-Guided Contrastive Learning for\\[0.4em]
    Clinically Aware Chest X-ray\\[0.4em]
    Image--Report Retrieval
  \par}

  \vspace{2cm}
  {\large \textbf{Diep Duc Lai}}\\[0.4em]
  {\normalsize Independent Researcher, Vietnam}

  \vspace{1.5cm}
  {\normalsize \today}

  \vspace{2.5cm}
  \rule{0.8\textwidth}{0.4pt}
  \vspace{0.8em}

  \begin{minipage}{0.85\textwidth}
    \begin{abstract}
      \noindent
      Cross-modal contrastive learning for chest X-ray image--report retrieval
      commonly suffers from the \emph{false-negative} problem: image--report
      pairs belonging to different patients but sharing the same pathological
      findings are treated as hard negatives, introducing incorrect gradient
      pressure and degrading clinically meaningful retrieval performance.
      We propose a \textbf{clustering-guided contrastive learning framework}
      to mitigate this issue on the IU-Xray dataset.
      Our model combines a \textbf{SwinV2-Base} image encoder with a
      \textbf{Bio\_ClinicalBERT} text encoder, projecting both modalities into
      a shared 512-dimensional embedding space through deep MLP projection
      heads. Multi-view study-level representations are formed via
      \emph{attention pooling} over frontal and lateral views.
      Training follows a three-phase curriculum:
      (i)~strict image--report alignment using multi-positive InfoNCE,
      (ii)~soft-target blending with IDF-weighted Jaccard pathology similarity,
      and (iii)~a clinical supervised contrastive loss with a gradual
      curriculum schedule that pulls same-disease, non-normal pairs together
      without compromising exact-pair retrieval.
      On a fixed patient-stratified split of IU-Xray (377 test studies),
      our model achieves a \textbf{clinical-valid mean R@1 of 59.86\%},
      cluster mean R@1 of 69.36\%, and strict mean R@1 of 3.85\%
      (MRR 9.01\%) on the test set.
      Qualitative analysis identifies \textbf{157 clinical rescue} cases
      in top-1 (64 image-to-text + 93 text-to-image), where the model
      retrieves a clinically related result from a different patient,
      directly validating the false-negative mitigation objective.

      \medskip
      \noindent\textbf{Keywords:}
      medical image--text retrieval,
      contrastive learning,
      false negatives,
      pathology clustering,
      chest X-ray,
      SwinV2,
      ClinicalBERT,
      curriculum learning.
    \end{abstract}
  \end{minipage}

  \vspace{0.8em}
  \rule{0.8\textwidth}{0.4pt}
\end{titlepage}

\tableofcontents
\newpage

%% ---------------------------------------------------------
%% 1. INTRODUCTION
%% ---------------------------------------------------------
\section{Introduction}
\label{sec:intro}

Automated medical image--text retrieval aims to associate diagnostic images
with their corresponding clinical reports in a shared embedding space.
Practical applications include case-based reasoning support, automated
report generation, and clinical decision support
systems~\citep{jing2018automatic}.

Multimodal contrastive learning in the style of
CLIP~\citep{radford2021learning} has emerged as the dominant paradigm for
this task. Methods such as ConVIRT~\citep{zhang2022contrastive} and
GLoRIA~\citep{huang2021gloria} train a model to bring together
\emph{same-patient} image--report pairs (positives) while pushing apart
\emph{different-patient} pairs (negatives).

\textbf{The False-Negative Problem.}
In clinical practice, patients from different studies frequently present
identical or highly similar pathological findings.
As illustrated in Figure~\ref{fig:false_negative}, standard InfoNCE treats
these \emph{clinically related} but \emph{patient-different} pairs as hard
negatives, penalizing the model for placing them nearby.
This produces incorrect gradient signals and suppresses the model's ability
to perform semantically meaningful, disease-level retrieval---the very
capability that is most valuable in a clinical context.

\begin{figure}[H]
  \centering
  \fbox{\parbox{0.88\textwidth}{
    \centering
    \vspace{0.4em}
    \textbf{False-Negative Problem in Medical Contrastive Learning}\\[0.6em]
    \small
    \begin{tabular}{p{3.5cm} c p{3.5cm}}
      \centering Patient A\\{\footnotesize Atelectasis} &
      \hspace{1cm} &
      \centering Patient B\\{\footnotesize Atelectasis}
    \end{tabular}\\[0.4em]
    \begin{tabular}{c}
      $\xrightarrow{\;\text{Standard InfoNCE}\;}$
      Treated as \textbf{negative} $\;\Rightarrow\;$ Pushed apart\\[0.3em]
      $\xrightarrow{\;\text{Proposed method}\;}$
      Treated as \textbf{soft positive} $\;\Rightarrow\;$ Not penalized
    \end{tabular}
    \vspace{0.4em}
  }}
  \caption{The false-negative problem in medical contrastive learning.
    Two patients sharing the same pathology (Atelectasis) are incorrectly
    penalized by standard InfoNCE. Our method assigns a soft positive target
    proportional to their pathological similarity.}
  \label{fig:false_negative}
\end{figure}

\textbf{Contributions.} This paper makes the following contributions:
\begin{enumerate}[label=(\roman*), leftmargin=2em]
  \item A \textbf{study-level multi-view retrieval pipeline} for IU-Xray
    using attention pooling over frontal and lateral X-ray views.
  \item A \textbf{hierarchical clustering loss} based on IDF-weighted Jaccard
    pathology similarity, generating soft training targets that reduce
    false-negative pressure from same-disease different-patient pairs.
  \item A \textbf{clinical supervised contrastive loss} with a curriculum
    schedule that gradually introduces and then decays disease-cluster
    supervision, preserving exact-pair retrieval ability.
  \item A \textbf{three-tier evaluation framework}---strict, cluster, and
    clinical-valid R@K---where \emph{clinical-valid} R@K is the primary
    metric directly measuring false-negative mitigation.
\end{enumerate}

The rest of this paper is organized as follows.
Section~\ref{sec:related} reviews related work.
Section~\ref{sec:method} describes the proposed method.
Section~\ref{sec:exp} presents experimental settings.
Section~\ref{sec:results} reports results.
Section~\ref{sec:discussion} discusses findings and limitations.
Section~\ref{sec:conclusion} concludes.

%% ---------------------------------------------------------
%% 2. RELATED WORK
%% ---------------------------------------------------------
\section{Related Work}
\label{sec:related}

\subsection{Multimodal Contrastive Learning}

CLIP~\citep{radford2021learning} established the foundation for
vision--language pretraining by optimizing InfoNCE~\citep{oord2018representation}
over 400 million web-crawled image--caption pairs.
ConVIRT~\citep{zhang2022contrastive} pioneered the application of contrastive
learning to radiology, aligning chest X-rays with clinical reports at the
global level.
GLoRIA~\citep{huang2021gloria} extended this with local attention-weighted
sub-regional alignment.
BiomedCLIP~\citep{zhang2023biomedclip} scaled medical vision--language
pretraining to 15 million figure--caption pairs from PubMed Central,
achieving strong zero-shot performance across multiple biomedical tasks.

\subsection{False Negatives in Contrastive Learning}

The false-negative problem has been studied in the general contrastive
learning literature.
\citet{khosla2020supervised} introduced Supervised Contrastive Learning
(SupCon), which uses class labels to identify multiple valid positives per
query in a batch.
\citet{chuang2020debiased} proposed a debiased estimator that corrects for
the probability of sampling true negatives.
\citet{robinson2021contrastive} studied hard negative sampling while
explicitly controlling false-negative contamination.

In radiology specifically, false negatives are particularly severe because:
(1) many patients share common pathologies (\eg{} cardiomegaly, atelectasis),
(2) the dataset is heavily label-imbalanced (IU-Xray: 57.88\% no-finding),
and (3) pathological semantics span both visual appearance and clinical
language simultaneously.
To our knowledge, no prior work has proposed an IDF-weighted Jaccard soft
target combined with a clinical curriculum for this problem on IU-Xray.

\subsection{Medical Visual--Linguistic Encoders}

\textbf{SwinV2}~\citep{liu2022swin} extends Swin Transformer with
post-normalization layer scaling and a continuous relative position bias,
enabling resolution transfer from low to high resolution (192$\to$384)
without degradation.
\textbf{Bio\_ClinicalBERT}~\citep{alsentzer2019publicly} is a BERT model
further pretrained on MIMIC-III clinical notes, making it well-suited for
the register of radiology reports.

\subsection{Chest X-ray Image--Report Retrieval}

The IU-Xray dataset~\citep{demner2016preparing} is the most widely used
benchmark for radiology report generation and retrieval.
Prior retrieval studies on IU-Xray typically report only strict R@K
(same-patient ground truth), which does not capture whether the model can
retrieve clinically related results from different patients.
This paper introduces \emph{clinical-valid} R@K to address this gap and
provide a more complete picture of retrieval quality.

%% ---------------------------------------------------------
%% 3. PROPOSED METHOD
%% ---------------------------------------------------------
\section{Proposed Method}
\label{sec:method}

Figure~\ref{fig:architecture} provides an overview of the model architecture
and training pipeline.

\begin{figure}[H]
  \centering
  \fbox{\parbox{0.93\textwidth}{
    \centering
    \vspace{0.5em}
    \textbf{Model Architecture Overview}\\[0.8em]
    \small
    \begin{tabular}{ccc}
      \makecell{\textbf{Image (Study)}\\ Frontal 384$\times$384\\ Lateral 384$\times$384} &
      $\;\xrightarrow{\;\text{SwinV2-Base}\;}\;$ &
      \makecell{$\vv{v}_f \in \mathbb{R}^{1024}$\\ $\vv{v}_l \in \mathbb{R}^{1024}$}
    \end{tabular}\\[0.3em]
    $\downarrow$ Projection MLP + View-type Embedding + Attention Pooling\\
    $\vv{z}_I \in \mathbb{R}^{512}$, \;$\norm{\vv{z}_I}_2 = 1$\\[0.7em]
    \begin{tabular}{ccc}
      \makecell{\textbf{Report} (max 128 tokens)} &
      $\;\xrightarrow{\;\text{Bio\_ClinicalBERT}\;}\;$ &
      \makecell{$\vv{u} \in \mathbb{R}^{768}$\\ (\texttt{[CLS]} token)}
    \end{tabular}\\[0.3em]
    $\downarrow$ Projection MLP\\
    $\vv{z}_T \in \mathbb{R}^{512}$, \;$\norm{\vv{z}_T}_2 = 1$\\[0.7em]
    Similarity: $\;s_{ij} = e^{\tau}\cdot\vv{z}_I^{(i)\top}\vv{z}_T^{(j)}$
    \quad (learnable $\tau$, $e^{\tau}\in[1,100]$)\\[0.5em]
    Loss: $\;\calL = \calL_{\text{base}} + w_c(e)\,\calL_{\text{clin}} + w_a\,\calL_{\text{aux}}$
    \vspace{0.4em}
  }}
  \caption{Architecture overview. Both image and text modalities are projected
    into a shared 512-d $\ell_2$-normalized embedding space.
    The loss combines a base contrastive term, a clinical supervised
    contrastive term with curriculum weight $w_c(e)$, and an auxiliary
    pathology prediction term.}
  \label{fig:architecture}
\end{figure}

\subsection{Data Representation}

Each sample in IU-Xray corresponds to one patient \emph{study}, comprising:
\begin{itemize}[leftmargin=1.5em]
  \item \textbf{Images:} up to one frontal and one lateral X-ray, stored as
    $384\times384$ RGB PNG after black-padding to preserve the full field of
    view without geometric distortion.
  \item \textbf{Report:} the concatenation of the \texttt{findings} and
    \texttt{impression} fields after text normalization (lowercasing, removal
    of placeholders, section headers, and short texts below 10 words).
  \item \textbf{Pathology vector:}
    $\vv{d}\in\{0,1\}^{12}$, derived from 14 CheXpert
    labels~\citep{irvin2019chexpert} via negation-aware keyword matching.
    ``No Finding'' is excluded; ``Enlarged Cardiomediastinum'' is merged into
    ``Cardiomegaly'' by taking the maximum.
\end{itemize}

A study is considered \textbf{normal} when $\sum_k d_k = 0$.

\subsection{Model Architecture}

\subsubsection{Image Encoder}

We use \textbf{SwinV2-Base}~\citep{liu2022swin},
specifically the checkpoint
\texttt{microsoft/swinv2-base-patch4-window12to24-192to384-22kto1k-ft},
pretrained on ImageNet-22K and fine-tuned at 384$\times$384 resolution.
The \texttt{pooler\_output} with dimension $\mathbb{R}^{1024}$ is used.
Gradient checkpointing is enabled to reduce GPU memory consumption.

\subsubsection{Text Encoder}

We use \textbf{Bio\_ClinicalBERT}~\citep{alsentzer2019publicly}
(\texttt{emilyalsentzer/Bio\_ClinicalBERT}), pretrained on MIMIC-III clinical
notes.
Reports are truncated to 128 tokens; the \texttt{[CLS]} embedding from the
final hidden layer ($\mathbb{R}^{768}$) is used as the report representation.

\subsubsection{Projection Head}

Both modalities pass through an identical two-layer MLP:
\begin{equation}
  \vv{z} = \frac{\text{MLP}(\vv{h})}{\norm{\text{MLP}(\vv{h})}_2},
  \label{eq:proj}
\end{equation}
where
\begin{equation}
  \text{MLP}(\vv{h}) =
    \vv{W}_2\,\text{Dropout}_{0.2}\!\bigl(
      \text{GELU}(\text{BN}(\vv{W}_1\vv{h}+\vv{b}_1))
    \bigr) + \vv{b}_2,
  \label{eq:mlp}
\end{equation}
with $\vv{W}_1\!\in\!\mathbb{R}^{1024\times d_{\mathrm{in}}}$,
$\vv{W}_2\!\in\!\mathbb{R}^{512\times 1024}$, BN = Batch Normalization.
The output $\vv{z}\!\in\!\mathbb{R}^{512}$ is $\ell_2$-normalized.

Similarity between an image--text pair is computed as:
\begin{equation}
  s_{ij} = e^{\tau} \cdot \vv{z}_I^{(i)\top} \vv{z}_T^{(j)},
  \label{eq:sim}
\end{equation}
where $\tau$ is a learnable scalar (\emph{logit\_scale}) clamped so that
$e^{\tau}\in[1,100]$, following~\citet{radford2021learning}.

\subsubsection{Multi-View Attention Pooling}

Each view (frontal / lateral) is independently encoded by SwinV2 and the
projection head to obtain $\vv{z}_k\in\mathbb{R}^{512}$.
A \emph{view-type embedding} $\vv{e}_k^{\mathrm{type}}\in\mathbb{R}^{512}$
(learned, one per view type: frontal, lateral, other, pad) is added.
The study embedding is computed by:
\begin{equation}
  \alpha_k =
    \frac{\exp\!\bigl(a(\vv{z}_k+\vv{e}_k^{\mathrm{type}})\bigr)}
         {\sum_{k'}\exp\!\bigl(a(\vv{z}_{k'}+\vv{e}_{k'}^{\mathrm{type}})\bigr)},
  \qquad
  \vv{z}_I = \sum_{k=1}^{K}\alpha_k\,\vv{z}_k,
  \label{eq:attn}
\end{equation}
where $a(\cdot) = \vv{w}_2^{\top}\,\text{GELU}(\vv{W}_1\,\cdot)$ is a
lightweight attention network ($\vv{W}_1\!\in\!\mathbb{R}^{256\times512}$,
$\vv{w}_2\!\in\!\mathbb{R}^{256}$), and $K\!\leq\!2$ is the number of valid
views (indicated by \texttt{view\_mask}).
Missing views are zero-padded and masked out.
This design allows the model to learn view importance adaptively rather
than relying on a fixed average.

\subsubsection{Auxiliary Pathology Heads}

Four lightweight MLP heads branch off the projected embeddings:
\texttt{img\_pathology\_head} and \texttt{img\_normal\_head} (from
$\vv{z}_I$) predict 12 pathology labels and a normal/abnormal binary label
respectively; \texttt{txt\_pathology\_head} and \texttt{txt\_normal\_head}
do the same from $\vv{z}_T$.
These heads enforce pathology-informative embeddings at a small auxiliary
weight ($w_a\!=\!0.02$), ensuring they do not dominate the retrieval
objective.

\subsection{Hierarchical Clustering Loss}
\label{sec:loss}

\subsubsection{Phase 1 --- Multi-Positive InfoNCE (Epochs 1--11)}

Before the clustering phase begins, the base loss is multi-positive InfoNCE:
\begin{equation}
  \calL_{\mathrm{InfoNCE}} =
    -\frac{1}{N}\sum_{i=1}^{N}
    \log\frac{\exp(s_{ii})}{\sum_{j=1}^{N}\exp(s_{ij})},
  \label{eq:infonce}
\end{equation}
computed bidirectionally (image$\to$text and text$\to$image).
Since $s_{ij}$ in~\eqref{eq:sim} already incorporates the learned
temperature, no additional temperature division is applied.
The ground-truth positive for query $i$ is the entry sharing the same
\texttt{patient\_id}.

\subsubsection{Phase 2 --- Hierarchical Clustering Loss (Epochs 12+)}
\label{sec:cluster_loss}

From epoch 12 onward, the hard binary ground-truth is replaced by a
\emph{soft target} that blends exact-patient matching with pathological
similarity:
\begin{equation}
  t_{ij} = (1-\alpha)\,\mathbb{1}[p_i=p_j]
           + \alpha\,S^{\mathrm{IDF}}_{ij},
  \qquad \alpha = 0.35.
  \label{eq:softtarget}
\end{equation}

The \textbf{IDF-weighted Jaccard similarity} between disease vectors
$\vv{d}_i,\vv{d}_j\in\{0,1\}^{12}$ is defined as:
\begin{equation}
  \boxed{
  S^{\mathrm{IDF}}(\vv{d}_i,\vv{d}_j) =
    \frac{\displaystyle\sum_{k=1}^{12}w_k\cdot\min(d_{ik},d_{jk})}
         {\displaystyle\sum_{k=1}^{12}w_k\cdot\max(d_{ik},d_{jk})},
  }
  \label{eq:jaccard}
\end{equation}
where the IDF weight for disease $k$ is:
\begin{equation}
  w_k = \frac{\log(1/f_k)}{\sum_{k'}\log(1/f_{k'})},
  \label{eq:idf}
\end{equation}
and $f_k$ is the prevalence of disease $k$ in the training set.
This weighting assigns higher importance to rare conditions
(\eg{} Pneumothorax $f\!=\!2.03\%$, Fracture $f\!=\!1.82\%$) over
more common ones (\eg{} Cardiomegaly $f\!=\!9.00\%$), reflecting the
greater clinical discriminative value of rare findings.

\textbf{Healthy cluster bonus.}
When both samples have $\sum_k d_k^{(i)} = \sum_k d_k^{(j)} = 0$ (both
normal), the Jaccard similarity is undefined.
We add a fixed similarity bonus:
\begin{equation}
  S^{\mathrm{IDF}}_{ij} \;\mathrel{+}=\; 0.4\cdot
    \mathbb{1}[\mathrm{normal}(i)]\cdot\mathbb{1}[\mathrm{normal}(j)],
  \label{eq:normalbonus}
\end{equation}
preventing the majority normal class (57.88\% of data) from scattering
in embedding space.

\subsubsection{Phase 3 --- Clinical Supervised Contrastive Loss (Epochs 15+)}
\label{sec:clinical_loss}

This is the core contribution targeting false-negative mitigation.

\textbf{Clinical mask.} For each pair $(i,j)$, the clinical positive
indicator is:
\begin{equation}
  M^{\mathrm{clin}}_{ij} =
    \underbrace{\mathbb{1}[\vv{d}_i\cdot\vv{d}_j > 0]}_{\text{disease overlap}}
    \cdot\;
    \mathbb{1}[i\neq j]
    \cdot\;
    \underbrace{\mathbb{1}[\mathrm{non{\text{-}}normal}(i)]
    \cdot\mathbb{1}[\mathrm{non{\text{-}}normal}(j)]}_{\text{both abnormal}}.
  \label{eq:clinmask}
\end{equation}
This selects pairs from \emph{different} studies that share at least one
disease label and are both abnormal.
Normal--normal pairs are excluded to avoid trivially inflating the metric.

\textbf{Loss.} For each non-normal query $i$ with at least one clinical
positive ($P_i=\{j:M^{\mathrm{clin}}_{ij}=1\}$, $|P_i|>0$):
\begin{equation}
  \calL_{\mathrm{clin}} =
    -\frac{1}{|\calQ|}\sum_{i\in\calQ}\frac{1}{|P_i|}
    \sum_{j\in P_i}\log\frac{\exp(s_{ij})}{\sum_{k=1}^{N}\exp(s_{ik})},
  \label{eq:clinloss}
\end{equation}
where $\calQ=\{i:|P_i|>0\}$.
The loss is computed bidirectionally.

\subsubsection{Auxiliary Pathology Loss}

\begin{equation}
  \calL_{\mathrm{aux}} =
    \frac{1}{4}\Bigl[
      \mathrm{BCE}(\hat{\vv{d}}_I,\vv{d})
      +\mathrm{BCE}(\hat{n}_I,n)
      +\mathrm{BCE}(\hat{\vv{d}}_T,\vv{d})
      +\mathrm{BCE}(\hat{n}_T,n)
    \Bigr]
    + 0.25\Bigl[
      \mathrm{MSE}(\sigma(\hat{\vv{d}}_I),\sigma(\hat{\vv{d}}_T))
      +\mathrm{MSE}(\sigma(\hat{n}_I),\sigma(\hat{n}_T))
    \Bigr],
  \label{eq:auxloss}
\end{equation}
where $\sigma$ is sigmoid.
The MSE term enforces cross-modal consistency in pathology predictions.

\textbf{Total loss:}
\begin{equation}
  \calL = \calL_{\mathrm{base}} + w_c(e)\,\calL_{\mathrm{clin}}
          + w_a\,\calL_{\mathrm{aux}},
  \quad w_a = 0.02.
  \label{eq:totalloss}
\end{equation}

\subsection{Clinical Curriculum Schedule}
\label{sec:curriculum}

Activating $\calL_{\mathrm{clin}}$ at full weight from the start risks
pulling the model toward broad disease clusters before it has learned
exact-pair alignment.
We therefore apply a ramp-up / hold / decay curriculum:
\begin{equation}
  w_c(e) =
  \begin{cases}
    0 & e < 15, \\[3pt]
    0.03 + 0.09\dfrac{e-15}{10} & 15 \leq e \leq 25, \\[6pt]
    0.12 & 26 \leq e \leq 60, \\[3pt]
    0.08 & 61 \leq e \leq 105, \\[3pt]
    0.04 & 106 \leq e \leq 150.
  \end{cases}
  \label{eq:curriculum}
\end{equation}

Table~\ref{tab:curriculum} summarizes the purpose of each phase.

\begin{table}[H]
  \centering
  \caption{Clinical curriculum schedule and training phase objectives.}
  \label{tab:curriculum}
  \renewcommand{\arraystretch}{1.1}
  \begin{tabular}{c L{4.5cm} c c}
    \toprule
    \textbf{Epochs} & \textbf{Primary objective} & \textbf{$w_c$} & \textbf{Phase} \\
    \midrule
    1--5   & Projection heads only (encoders frozen) & 0.00 & Warm-up \\
    6--11  & Multi-positive InfoNCE, all unfrozen     & 0.00 & Strict align \\
    12--14 & Soft clustering targets activated        & 0.00 & Cluster init \\
    15--25 & Clinical loss ramping                   & 0.03$\to$0.12 & Clinical ramp \\
    26--60 & Peak clinical supervision               & 0.12 & Clinical peak \\
    61--105 & Clinical weight decay                 & 0.08 & Decay \\
    106--150 & Final strict recovery                 & 0.04 & Fine-tune \\
    \bottomrule
  \end{tabular}
\end{table}

The rationale is validated by our results (Section~\ref{sec:results}):
strict R@1 reaches its validation peak at epoch 120, \emph{after} the
clinical decay begins at epoch 60, confirming that reducing $w_c$ late in
training allows the model to recover precise image--report alignment.

%% ---------------------------------------------------------
%% 4. EXPERIMENTS
%% ---------------------------------------------------------
\section{Experiments}
\label{sec:exp}

\subsection{Dataset}

The \textbf{IU-Xray} dataset~\citep{demner2016preparing}
(Indiana University Chest X-ray Collection) contains chest radiographs
collected from Indiana University hospital network, paired with
de-identified radiology reports.

\textbf{Preprocessing.}
Images are converted to grayscale, then to RGB, padded to a square aspect
ratio with a black background, and saved at $384\!\times\!384$ pixels to
match SwinV2-Base's native resolution without cropping.
Reports are lower-cased, stripped of section headers
(\eg{} ``findings:'', ``impression:'') and age/gender tokens, and filtered
to retain only reports with $\geq\!10$ words.
Pathology labels are generated by negation-aware keyword matching over 14
CheXpert categories and then reduced to a 12-dimensional disease vector.

Dataset statistics are summarized in Table~\ref{tab:dataset}.

\begin{table}[H]
  \centering
  \caption{IU-Xray dataset statistics after preprocessing.}
  \label{tab:dataset}
  \renewcommand{\arraystretch}{1.1}
  \begin{tabular}{lr}
    \toprule
    \textbf{Statistic} & \textbf{Value} \\
    \midrule
    Total images (384$\times$384 PNG) & 7,322 \\
    Unique patient / study IDs        & 3,772 \\
    Unique reports                    & 3,018 \\
    Words per report (min / mean / max) & 10 / 37.27 / 223 \\
    Frontal views                     & 3,738 \\
    Lateral views                     & 3,584 \\
    \midrule
    \multicolumn{2}{l}{\textit{Patient-stratified split (seed = 42)}} \\
    Train (study-level)  & 3,018 \\
    Validation           & 377   \\
    Test                 & 377   \\
    \bottomrule
  \end{tabular}
\end{table}

Pathology label prevalence is shown in Table~\ref{tab:labels}.

\begin{table}[H]
  \centering
  \caption{Prevalence of 14 CheXpert pathology labels in IU-Xray (7,322 images).}
  \label{tab:labels}
  \renewcommand{\arraystretch}{1.05}
  \begin{tabular}{lrr}
    \toprule
    \textbf{Label} & \textbf{Count} & \textbf{\%} \\
    \midrule
    No Finding                  & 4,238 & 57.88 \\
    Cardiomegaly                &   659 &  9.00 \\
    Lung Lesion                 & 1,014 & 13.85 \\
    Lung Opacity                &   942 & 12.87 \\
    Atelectasis                 &   732 & 10.00 \\
    Consolidation               &   402 &  5.49 \\
    Pleural Effusion            &   415 &  5.67 \\
    Support Devices             &   224 &  3.06 \\
    Edema                       &   185 &  2.53 \\
    Pneumonia                   &   164 &  2.24 \\
    Pneumothorax                &   149 &  2.03 \\
    Fracture                    &   133 &  1.82 \\
    Pleural Other               &    65 &  0.89 \\
    Enlarged Cardiomediastinum  &    17 &  0.23 \\
    \bottomrule
  \end{tabular}
\end{table}

\subsection{Implementation Details}

\textbf{Hardware.} NVIDIA A100 (40 GB VRAM), PyTorch with
\texttt{torch.amp.autocast} + \texttt{GradScaler} mixed-precision training.

\textbf{Hyperparameters} are listed in Table~\ref{tab:hparams}.

\begin{table}[H]
  \centering
  \caption{Training hyperparameters for the v8 run
    (\texttt{v8\_curriculum\_a100\_150ep\_seed42\_bs32\_v1}).}
  \label{tab:hparams}
  \renewcommand{\arraystretch}{1.1}
  \begin{tabular}{ll}
    \toprule
    \textbf{Hyperparameter} & \textbf{Value} \\
    \midrule
    Epochs                          & 150 \\
    Batch size                      & 32 \\
    Gradient accumulation steps     & 4 \\
    Effective batch size            & 128 \\
    Backbone freeze epochs          & 5 \\
    Evaluation frequency            & every 5 epochs \\
    Random seed                     & 42 \\
    \midrule
    LR --- image encoder            & $1\times10^{-5}$ \\
    LR --- text encoder             & $1\times10^{-4}$ \\
    LR --- projection heads         & $2\times10^{-4}$ \\
    Weight decay                    & 0.01 \\
    Max gradient norm               & 1.0 \\
    Optimizer                       & AdamW \\
    LR scheduler (post-unfreeze)    & Linear warmup (5\%) + CosineAnnealing \\
    $\eta_{\min}$                   & $10^{-8}$ \\
    \midrule
    $\alpha$ (soft-target blend)    & 0.35 \\
    Cluster temperature             & 1.0 \\
    $w_a$ (auxiliary weight)        & 0.02 \\
    \midrule
    Embedding dimension             & 512 \\
    Projection hidden dimension     & 1,024 \\
    Projection dropout              & 0.2 \\
    \bottomrule
  \end{tabular}
\end{table}

\textbf{Checkpoint strategy.}
Three checkpoints are maintained:
\texttt{best.pt} (best validation strict R@1),
\texttt{best\_balanced.pt} (best strict $+\,0.1\times$clinical-valid,
when strict $\geq4.0\%$), and
\texttt{best\_clinical\_valid.pt} (best clinical-valid, when
strict $\geq3.5\%$).
The primary reported checkpoint is \texttt{best.pt} (epoch 120).

\textbf{Disabled components.}
For transparency, we note that the codebase contains
\emph{DualPrototypeBank}, \emph{IntraClusterRankingLoss}, and
\emph{Hard Negative Mining} components, but all are \emph{disabled} in the
v8 final run
(\texttt{proto\_start=999}, \texttt{rank\_start=999},
\texttt{mine\_start=999}).
The results reported here do \emph{not} rely on these components; they are
reserved for future ablation studies.

\subsection{Evaluation Metrics}
\label{sec:metrics}

All evaluations are at study-level: each patient is one query, and the
query embedding is the attention-pooled image embedding.
We report metrics in both directions: image-to-text (i2t) and
text-to-image (t2i).

\textbf{M1 --- Strict Retrieval.}
$G^{\mathrm{strict}}_{ij} = \mathbb{1}[p_i=p_j]$.
Only the same-patient study is a ground-truth hit.
Metrics: R@1, R@5, R@10, MRR.

\textbf{M2 --- Cluster Retrieval.}
$G^{\mathrm{clus}}_{ij} = \mathbb{1}[\vv{d}_i\!\cdot\!\vv{d}_j>0]
  + \mathbb{1}[\mathrm{normal}(i)\wedge\mathrm{normal}(j)]$.
Both disease-overlapping pairs and normal--normal pairs are hits.

\textbf{M3 --- Clinical-Valid Retrieval (primary metric).}
$G^{\mathrm{clin}}_{ij} = \mathbb{1}[\vv{d}_i\!\cdot\!\vv{d}_j>0]
  \wedge\mathbb{1}[\mathrm{non{\text{-}}normal}(i)]
  \wedge\mathbb{1}[\mathrm{non{\text{-}}normal}(j)]$.
Computed only for non-normal queries with $\geq\!1$ clinical positive
in the test set ($n\!=\!147$ per direction).
This metric directly measures false-negative mitigation.

\textbf{M4 --- Clinical-All Retrieval.}
Same mask as M3 but evaluated over all 377 queries (including those with
no clinical positive). Used as a secondary metric.

\textbf{Qualitative bucket analysis.}
Each top-1 result is classified into one of four mutually exclusive buckets:
\emph{strict hit} (same patient), \emph{clinical rescue} (different patient,
clinical mask satisfied), \emph{cluster rescue} (cluster mask satisfied
but not clinical), or \emph{miss}.

%% ---------------------------------------------------------
%% 5. RESULTS
%% ---------------------------------------------------------
\section{Results}
\label{sec:results}

\subsection{Quantitative Results on the Test Set}

Table~\ref{tab:main} reports the full test-set metrics for the primary
checkpoint \texttt{best.pt} (epoch 120, 377 test studies).

\begin{table}[H]
  \centering
  \caption{Test-set retrieval performance of the proposed model (v8,
    \texttt{best.pt}, epoch 120). \textbf{Mean} = average of i2t and t2i.
    The primary metric is \emph{clinical-valid} (shaded).}
  \label{tab:main}
  \renewcommand{\arraystretch}{1.15}
  \setlength{\tabcolsep}{5pt}
  \begin{tabular}{l ccc ccc >{\columncolor{gray!12}}c >{\columncolor{gray!12}}c >{\columncolor{gray!12}}c}
    \toprule
    & \multicolumn{3}{c}{\textbf{Strict}} &
      \multicolumn{3}{c}{\textbf{Cluster}} &
      \multicolumn{3}{c}{\textbf{Clinical-Valid}} \\
    \cmidrule(lr){2-4}\cmidrule(lr){5-7}\cmidrule(lr){8-10}
    & R@1 & R@5 & R@10 & R@1 & R@5 & R@10 & R@1 & R@5 & R@10 \\
    \midrule
    i2t & 3.45 & 10.08 & 17.77 & 68.17 & 79.58 & 85.15 & 49.66 & 60.54 & 69.39 \\
    t2i & 4.24 & 12.20 & 18.83 & 70.56 & 97.08 & 99.20 & 70.07 & 95.24 & 97.96 \\
    \midrule
    \textbf{Mean}
      & \textbf{3.85} & \textbf{11.14} & \textbf{18.30}
      & \textbf{69.36} & \textbf{88.33} & \textbf{92.18}
      & \textbf{59.86} & \textbf{77.89} & \textbf{83.67} \\
    \bottomrule
  \end{tabular}
  \vspace{0.3em}\\
  \small
  \textit{Strict MRR = 9.01\%;\quad
  Clinical-all mean R@1 = 23.34\%;\quad
  Clinical-valid: 147\,/\,377 queries per direction have $\geq$1 positive.}
\end{table}

\textbf{Best validation metrics} (across all epochs):
strict mean R@1 = \textbf{4.91\%} (epoch 120);
clinical-valid mean R@1 = \textbf{52.11\%} (epoch 90);
cluster mean R@1 = \textbf{65.25\%} (epoch 40).

\subsection{Comparison with Baseline}

Table~\ref{tab:comparison} compares the proposed model against v7, which
uses the same architecture and split but a less refined curriculum schedule.

\begin{table}[H]
  \centering
  \caption{Comparison of v8 (proposed) vs.\ v7 (baseline) on the test set.
    Both models use the same patient-stratified split (seed=42).}
  \label{tab:comparison}
  \renewcommand{\arraystretch}{1.1}
  \begin{tabular}{lcccc}
    \toprule
    \textbf{Model} &
    \textbf{Strict Mean R@1} &
    \textbf{Cluster Mean R@1} &
    \textbf{Clin.-Valid Mean R@1} &
    \textbf{MRR} \\
    \midrule
    V7 (best-balanced ckpt) & 3.32\% & 67.37\% & 59.86\% & --- \\
    \textbf{V8 (proposed, ep.~120)} & \textbf{3.85\%} & \textbf{69.36\%} & \textbf{59.86\%} & \textbf{9.01\%} \\
    \midrule
    $\Delta$ & \textbf{+0.53 pp} & \textbf{+1.99 pp} & 0.00 pp & --- \\
    \bottomrule
  \end{tabular}
\end{table}

V8 improves strict mean R@1 by \textbf{+0.53 pp} and cluster mean R@1 by
\textbf{+1.99 pp} while maintaining clinical-valid mean R@1 at the same
high level of 59.86\%.
This demonstrates that the curriculum decay in v8 successfully preserves
and improves exact retrieval without sacrificing clinical-semantic retrieval.

\subsection{Qualitative Bucket Analysis}

Table~\ref{tab:bucket} classifies every top-1 retrieval result into the
four quality buckets described in Section~\ref{sec:metrics}.

\begin{table}[H]
  \centering
  \caption{Top-1 qualitative bucket analysis over all 377 test studies.
    \emph{Clinical rescue} cases are the key validation of the
    false-negative mitigation objective.}
  \label{tab:bucket}
  \renewcommand{\arraystretch}{1.1}
  \begin{tabular}{lcccc}
    \toprule
    \textbf{Direction} &
    \textbf{Strict Hit} &
    \textbf{Clinical Rescue} &
    \textbf{Cluster Rescue} &
    \textbf{Miss} \\
    \midrule
    Image $\to$ Text & 13 (3.4\%)  & \textbf{64 (17.0\%)}  & 180 (47.7\%) & 120 (31.8\%) \\
    Text $\to$ Image & 16 (4.2\%)  & \textbf{93 (24.7\%)} & 157 (41.6\%) & 111 (29.4\%) \\
    \midrule
    \textbf{Combined} & 29 (3.8\%) & \textbf{157 (20.8\%)} & 337 (44.7\%) & 231 (30.7\%) \\
    \bottomrule
  \end{tabular}
\end{table}

Of 754 total top-1 retrievals, \textbf{157 (20.8\%) are clinical rescues}:
the model retrieved the wrong patient but a clinically related study.
An additional 337 (44.7\%) are cluster rescues.
Together, 523\,/\,754 top-1 results (69.4\%) are at least clinically
reasonable, with only 231 complete misses.

\subsection{Training Dynamics}

Key milestones observed in the validation curves:

\begin{itemize}[leftmargin=1.5em]
  \item \textbf{Epoch 12:} clustering soft targets activated; cluster R@1
    begins rising.
  \item \textbf{Epoch 15--25:} clinical weight ramps up; clinical-valid R@1
    improves steadily.
  \item \textbf{Epoch 40:} cluster R@1 reaches its validation peak
    (65.25\%).
  \item \textbf{Epoch 60:} clinical weight begins decay; strict R@1 starts
    recovering.
  \item \textbf{Epoch 90:} clinical-valid R@1 peaks on validation (52.11\%).
  \item \textbf{Epoch 120:} \textbf{strict R@1 peaks on validation (4.91\%)}
    --- \texttt{best.pt} checkpoint saved.
\end{itemize}

The late strict peak at epoch 120, after the clinical weight has been
decaying for 60 epochs, directly validates the curriculum design rationale.

%% ---------------------------------------------------------
%% 6. DISCUSSION
%% ---------------------------------------------------------
\section{Discussion}
\label{sec:discussion}

\subsection{Interpretation of Clinical-Valid R@1}

A clinical-valid mean R@1 of 59.86\% over 147 queries per direction means
that in more than half of cases where a clinically related result exists in
the test set, the model retrieves it at rank 1.
In a clinical workflow, retrieving a disease-matched case from a different
patient is often more useful than retrieving the exact same study---a
radiologist needs a visually and clinically similar reference, not a
historical duplicate.

\subsection{Why Strict R@1 Remains Low}

Strict R@1 = 3.85\% is modest, attributable to: (1) the small scale of
IU-Xray ($\sim$3,000 training studies) relative to the task complexity;
(2) high textual homogeneity---many normal reports contain nearly identical
language, making exact discrimination difficult; and (3) the inherent
trade-off between learning exact-pair alignment and disease-cluster
semantics.
We report this honestly as a known limitation rather than masking it with
cluster or clinical-valid numbers alone.

\subsection{Role of IDF Weighting and the Healthy Cluster}

IDF weighting ensures that rare diseases (\eg{} Pneumothorax 2.03\%,
Fracture 1.82\%) contribute proportionally more to the similarity score
than common diseases (\eg{} Cardiomegaly 9.00\%).
This is clinically motivated: shared presence of a rare finding is a much
stronger indicator of clinical relatedness than shared cardiomegaly.
The healthy cluster bonus (Eq.~\ref{eq:normalbonus}) prevents the dominant
normal class (57.88\%) from fragmenting in embedding space, which would
otherwise inflate cluster R@1 while hurting the discriminative representations.

\subsection{Curriculum vs.\ Fixed Clinical Weight}

The curriculum schedule (ramp $\to$ hold $\to$ decay) is critical.
A fixed high $w_c$ from training onset would pull the model toward coarse
disease clusters before exact alignment has been established, degrading
strict R@1 permanently.
The epoch-120 strict peak occurring \emph{after} 60 epochs of clinical
decay empirically confirms that late-phase reduction of $w_c$ allows the
model to re-sharpen exact-pair discrimination.

\subsection{Limitations}

\begin{enumerate}[leftmargin=2em, label=\arabic*.]
  \item \textbf{Dataset scale.}
    IU-Xray contains only $\sim$3,000 training studies---far smaller than
    large-scale medical CLIP datasets (BiomedCLIP: 15M pairs).
    Performance may improve substantially with more data.
  \item \textbf{Weak pathology supervision.}
    Disease labels are generated by rule-based keyword matching, not
    physician annotation.
    Complex negation patterns may yield label noise.
  \item \textbf{No published-baseline comparison.}
    V8 is compared only against an internal v7 baseline.
    Direct comparison with ConVIRT, GLoRIA, or BiomedCLIP on the same
    split is left for future work.
  \item \textbf{Single-dataset evaluation.}
    Generalization to larger datasets such as MIMIC-CXR has not been tested.
  \item \textbf{Cluster metric inflation.}
    Cluster R@1 partially benefits from the large normal class;
    clinical-valid R@1 should be treated as the primary claim.
  \item \textbf{Disabled components.}
    Prototype bank and hard negative mining are in the codebase but not
    activated in v8; ablation studies for these are deferred to future work.
\end{enumerate}

%% ---------------------------------------------------------
%% 7. CONCLUSION
%% ---------------------------------------------------------
\section{Conclusion}
\label{sec:conclusion}

We have presented a \textbf{clustering-guided contrastive learning framework}
for chest X-ray image--report retrieval that directly addresses the
false-negative problem inherent in standard contrastive training.
The model combines a SwinV2-Base image encoder with Bio\_ClinicalBERT,
integrates multi-view fusion via attention pooling, and is trained with a
three-phase curriculum that introduces and then decays clinical supervision.

On a fixed patient-stratified split of IU-Xray:
\begin{itemize}[leftmargin=1.5em]
  \item \textbf{Clinical-valid mean R@1 = 59.86\%} (primary metric, 147 queries/direction).
  \item \textbf{Cluster mean R@1 = 69.36\%}.
  \item \textbf{Strict mean R@1 = 3.85\%}; \textbf{MRR = 9.01\%}.
  \item \textbf{157 clinical rescue cases} in top-1 (64 i2t + 93 t2i / 377 test studies).
\end{itemize}

The results validate the thesis that clustering-guided supervision
substantially improves disease-level semantic retrieval while the
curriculum preserves exact-pair alignment.
Promising directions for future work include: systematic ablation of each
proposed component (IDF Jaccard weighting, healthy cluster bonus, curriculum
schedule), activation and evaluation of the prototype bank and hard negative
mining modules, and cross-dataset generalization experiments on MIMIC-CXR.

\section*{Acknowledgements}

The author thanks the thesis advisor and research group for their guidance.
Experiments were conducted on NVIDIA A100 GPU through Google Colaboratory
and Kaggle platforms.

%% ---------------------------------------------------------
%% REFERENCES
%% ---------------------------------------------------------
\begin{thebibliography}{99}

\bibitem[Alsentzer et~al.(2019)]{alsentzer2019publicly}
Alsentzer, E., Murphy, J.~R., Boag, W., Weng, W.-H., Jin, D., Naumann, T.,
  and McDermott, M. (2019).
\newblock Publicly available clinical {BERT} embeddings.
\newblock In \textit{Proc.\ 2nd Clinical NLP Workshop (ACL)}, pp.\ 72--78.

\bibitem[Chuang et~al.(2020)]{chuang2020debiased}
Chuang, C.-Y., Robinson, J., Lin, Y.-C., Torralba, A., and Jegelka, S.
  (2020).
\newblock Debiased contrastive learning.
\newblock In \textit{Advances in Neural Information Processing Systems
  (NeurIPS)}, Vol.~33, pp.\ 8765--8775.

\bibitem[Demner-Fushman et~al.(2016)]{demner2016preparing}
Demner-Fushman, D., Kohli, M.~D., Rosenman, M.~B., Shooshan, S.~E.,
  Rodriguez, L., Antani, S., Thoma, G.~R., and McDonald, C.~J. (2016).
\newblock Preparing a collection of radiology examinations for distribution
  and retrieval.
\newblock \textit{J.\ Am.\ Med.\ Inform.\ Assoc.\ (JAMIA)}, 23(2):304--310.

\bibitem[Huang et~al.(2021)]{huang2021gloria}
Huang, S.-C., Shen, L., Lungren, M.~P., and Yeung, S. (2021).
\newblock {GloRIA}: A multimodal global-local representation learning
  framework for label-efficient medical image recognition.
\newblock In \textit{Proc.\ IEEE/CVF ICCV}, pp.\ 3942--3951.

\bibitem[Irvin et~al.(2019)]{irvin2019chexpert}
Irvin, J., Rajpurkar, P., Ko, M., Yu, Y., et~al. (2019).
\newblock {CheXpert}: A large chest radiograph dataset with uncertainty
  labels and expert comparison.
\newblock In \textit{Proc.\ AAAI Conference on AI}, Vol.~33, pp.\ 590--597.

\bibitem[Jing et~al.(2018)]{jing2018automatic}
Jing, B., Xie, P., and Xing, E. (2018).
\newblock On the automatic generation of medical imaging reports.
\newblock In \textit{Proc.\ ACL 2018}, pp.\ 2577--2586.

\bibitem[Khosla et~al.(2020)]{khosla2020supervised}
Khosla, P., Tian, Y., Wang, X., Liu, C., Isola, P., Krishnamurthy, A.,
  and Tian, Y. (2020).
\newblock Supervised contrastive learning.
\newblock In \textit{Advances in Neural Information Processing Systems
  (NeurIPS)}, Vol.~33, pp.\ 18661--18673.

\bibitem[Liu et~al.(2022)]{liu2022swin}
Liu, Z., Hu, H., Lin, Y., Yao, Z., Xie, Z., Wei, Y., Ning, J., Cao, Y.,
  Zhang, Z., Dong, L., et~al. (2022).
\newblock Swin transformer {V2}: Scaling up capacity and resolution.
\newblock In \textit{Proc.\ IEEE/CVF CVPR}, pp.\ 12009--12019.

\bibitem[Oord et~al.(2018)]{oord2018representation}
Oord, A.~v.~d., Li, Y., and Vinyals, O. (2018).
\newblock Representation learning with contrastive predictive coding.
\newblock \textit{arXiv preprint arXiv:1807.03748}.

\bibitem[Radford et~al.(2021)]{radford2021learning}
Radford, A., Kim, J.~W., Hallacy, C., Ramesh, A., Goh, G., Agarwal, S.,
  Sastry, G., Askell, A., Mishkin, P., Clark, J., et~al. (2021).
\newblock Learning transferable visual models from natural language
  supervision.
\newblock In \textit{Proc.\ ICML}, pp.\ 8748--8763.

\bibitem[Robinson et~al.(2021)]{robinson2021contrastive}
Robinson, J., Chuang, C.-Y., Sra, S., and Jegelka, S. (2021).
\newblock Contrastive learning with hard negative samples.
\newblock In \textit{Proc.\ ICLR 2021}.

\bibitem[Zhang et~al.(2022)]{zhang2022contrastive}
Zhang, Y., Jiang, H., Miura, Y., Manning, C.~D., and Langlotz, C.~P.
  (2022).
\newblock Contrastive learning of medical visual representations from
  paired images and text.
\newblock In \textit{Proc.\ Machine Learning for Healthcare (MLHC)},
  pp.\ 2--25.

\bibitem[Zhang et~al.(2023)]{zhang2023biomedclip}
Zhang, S., Xu, Y., Usuyama, N., Bagga, J.~K., Tinn, R., Preston, S.,
  Rao, R., Wei, M., Valluri, N., Wong, C., et~al. (2023).
\newblock Large-scale domain-specific pretraining for biomedical
  vision--language processing.
\newblock \textit{arXiv preprint arXiv:2303.00915}.

\end{thebibliography}

%% ---------------------------------------------------------
%% APPENDIX
%% ---------------------------------------------------------
\appendix

\section{Full Hyperparameter Configuration (v8)}
\label{app:config}

\begin{table}[H]
  \centering
  \caption{Complete hyperparameter configuration of the primary v8 run.}
  \label{tab:fullconfig}
  \small\renewcommand{\arraystretch}{1.05}
  \begin{tabular}{ll}
    \toprule
    \textbf{Parameter} & \textbf{Value} \\
    \midrule
    \multicolumn{2}{l}{\textit{General}} \\
    epochs / batch\_size / grad\_accum   & 150 / 32 / 4 \\
    effective batch size / seed          & 128 / 42 \\
    freeze\_ep / eval\_every             & 5 / 5 \\
    \midrule
    \multicolumn{2}{l}{\textit{Optimizer}} \\
    lr\_head / lr\_enc / lr\_text        & 2e-4 / 1e-5 / 1e-4 \\
    weight\_decay / max\_grad\_norm      & 0.01 / 1.0 \\
    \midrule
    \multicolumn{2}{l}{\textit{Clustering}} \\
    cluster\_start / CLUSTER\_ALPHA      & 12 / 0.35 \\
    proto\_start (disabled)              & 999 \\
    rank\_start / rank\_weight (disabled)& 999 / 0.0 \\
    mine\_start / hn\_frac (disabled)    & 999 / 0.0 \\
    \midrule
    \multicolumn{2}{l}{\textit{Clinical curriculum}} \\
    clinical\_start                      & 15 \\
    ramp start weight / peak weight      & 0.03 / 0.12 \\
    ramp end / hold end                  & 25 / 60 \\
    mid weight / decay end               & 0.08 / 105 \\
    final weight                         & 0.04 \\
    \midrule
    \multicolumn{2}{l}{\textit{Loss weights}} \\
    aux\_weight                          & 0.02 \\
    \bottomrule
  \end{tabular}
\end{table}

\section{IDF Weights for Pathology Labels}
\label{app:idf}

\begin{table}[H]
  \centering
  \caption{Relative IDF weights of the 12 pathology dimensions in the
    disease vector, derived from training-set prevalence
    (Table~\ref{tab:labels}).
    Rare diseases receive higher weight, reflecting their greater
    discriminative clinical value.}
  \label{tab:idfweights}
  \small\renewcommand{\arraystretch}{1.05}
  \begin{tabular}{lrl}
    \toprule
    \textbf{Pathology} & \textbf{Prevalence (\%)} & \textbf{Relative IDF weight} \\
    \midrule
    Pleural Other          & 0.89  & Highest  $\blacksquare\blacksquare\blacksquare\blacksquare\blacksquare$ \\
    Fracture               & 1.82  & Very high $\blacksquare\blacksquare\blacksquare\blacksquare$ \\
    Pneumothorax           & 2.03  & Very high $\blacksquare\blacksquare\blacksquare\blacksquare$ \\
    Pneumonia              & 2.24  & High $\blacksquare\blacksquare\blacksquare$ \\
    Edema                  & 2.53  & High $\blacksquare\blacksquare\blacksquare$ \\
    Support Devices        & 3.06  & Medium $\blacksquare\blacksquare$ \\
    Consolidation          & 5.49  & Medium $\blacksquare\blacksquare$ \\
    Pleural Effusion       & 5.67  & Medium $\blacksquare\blacksquare$ \\
    Cardiomegaly           & 9.00  & Low $\blacksquare$ \\
    Atelectasis            & 10.00 & Low $\blacksquare$ \\
    Lung Opacity           & 12.87 & Low $\blacksquare$ \\
    Lung Lesion            & 13.85 & Lowest $\square$ \\
    \bottomrule
  \end{tabular}
\end{table}

\section{Checkpoint Comparison on Test Set}
\label{app:ckpts}

\begin{table}[H]
  \centering
  \caption{Comparison of all three saved checkpoints on the test set.
    \texttt{best.pt} is the primary checkpoint; \texttt{best\_balanced.pt}
    coincides with it at epoch 120.}
  \label{tab:ckpts}
  \renewcommand{\arraystretch}{1.1}
  \begin{tabular}{lcccc}
    \toprule
    \textbf{Checkpoint} & \textbf{Epoch} &
    \textbf{Strict} & \textbf{Cluster} & \textbf{Clin.-Valid} \\
    & & \textbf{Mean R@1} & \textbf{Mean R@1} & \textbf{Mean R@1} \\
    \midrule
    \texttt{best.pt}                  & 120 & \textbf{3.85\%} & \textbf{69.36\%} & 59.86\% \\
    \texttt{best\_balanced.pt}        & 120 & 3.85\% & 69.36\% & 59.86\% \\
    \texttt{best\_clinical\_valid.pt} & 90  & 2.39\% & 65.78\% & 59.52\% \\
    \bottomrule
  \end{tabular}
\end{table}

\end{document}
